\section{Extension to Schreier Graphs}
In this section, we extend the framework from the Boolean hypercube $\Omega=\{\pm 1\}^n$ to Schreier graphs. This more general setting was also considered in~\cite{ODonnellWimmer13,ODonnellWimmerSharp13,CorderoErausquinLedoux12}, which established a KKL theorem for Schreier graphs given a log-Sobolev constant. Here, we derive an Eldan--Gross inequality in this setting under a slightly stronger assumption.

\subsection{Basic definitions}
Let $G$ be a group acting on a finite set $X$, and $U$ be a generating set for $G$ that is symmetric, i.e., $U^{-1} = U$. The \emph{Schreier graph} $\Sch(G,X,U)$ is a graph with vertex set $V = X$ and edge set $E = \{(x,gx) \mid x \in X, g \in U\}$. Suppose that $U = \bigcup_{i=1}^n U_i$ where each $U_i$ is symmetric, i.e., $U_i^{-1} = U_i$. We remark that $\{U_i\}$ is not necessarily a partition of $U$. {\color{red} Though, each group element $u \in U$ will appears exactly in $r$ subsets $U_i$ for some $r \ge 1$.}

Schreier graphs encompass a broad range of discrete structures, including product spaces such as $[q]^d$ and slices of the Boolean hypercube such as $\binom{[n]}{k}$.
\begin{example}\label{example:SGraph}
    Let $G = X = \Z_m^n$ and $U_i = \{k e_i: k \in [m-1]\}$. The Schreier graph $\Sch(G,X,U)$ is a Hamming graph.
\end{example}

\begin{example}
    Let $G = S_n$, $X = \binom{[n]}{k}$ and $U_i = \{(ij) \mid j \in [n] \setminus \{i\}\}$. The Schreier graph $\Sch(G,X,U)$ is the exchange graph on the Boolean slice $\binom{[n]}{k}$, and the corresponding random walk is the standard exchange walk.
\end{example}

Similar to the Boolean hypercube case, we define the carr\'e du champ operators, and influences as well as surface area of Boolean functions for the Schreier graphs.

\begin{definition}[$\Gamma$-calculus]
Let $\Sch(G,X,U)$ be a Schreier graph with $U = \bigcup_{i=1}^n U_i$. Given the transition rates $\{q_u: \Omega \to \mathbb{R}_{+}\}_{u \in U}$ and weights $\{w_u: \Omega \to \mathbb{R}_+\}_{u \in U}$, we define the operator $L_w$ and carr\'e du champ operators as follows:
\begin{itemize}
    \item The operator $L_w$ on functions $f: X \to \mathbb{R}$ satisfies
    \begin{align*}
        L_w f(x) = \sum_{u \in U} w_u(x) q_u(x) (f(x^u) - f(x)),
    \end{align*}
    where $x^u = ux$, and we omit the subscript $w$ if $w \equiv 1$.
    \item For any $f,g : X \to \mathbb{R}$, the carr\'e du champ operator $\Gamma_w$ is defined as
    \begin{align*}
        \Gamma_w (f,g) = \frac{1}{2}\tp{L_w (fg) - f L_w g - g L_w f}.
    \end{align*}
    By a straightforward calculation,
    \begin{align*}
        \Gamma_w(f,g)(x) = \frac{1}{2r} \sum_{i \in [n]} \sum_{u \in U_i} w_u(x) q_u(x) \tp{f(x^u) - f(x)}\tp{g(x^u)-g(x)}.
    \end{align*}
    Hence, the operator $\Gamma_{w,i}$ for the $i$-th ``coordinate'' is defined as
    \begin{align*}
    \Gamma_{w,i} (f,g) (x) = \frac{1}{2 r}\sum_{u \in U_i} w_u(x) q_u(x) \tp{f(x^u) - f(x)}\tp{g(x^u)-g(x)}.
    \end{align*}
\end{itemize}
\end{definition}

Similar to~\Cref{def:weighted-gradient}, we define the absolute weighted discrete gradient.
\begin{definition}
    For any function $f:X \to \R_+$, the absolute weighted discrete gradient is given by
    \begin{align*}
        \abs{\nabla_w} f = (\abs{\delta_{w,i}} f)_{i \in [n]}, \quad \text{where $\abs{\delta_{w,i}} f = \sqrt{\Gamma_{w,i}(f)}$}.
    \end{align*}
\end{definition}

With above definitions, we are now ready to define the influences and surface area.


\begin{definition}
    Let $f: X \to \mathbb{R}$ be a function.

    \begin{enumerate}[(1)]
        \item The \textit{total influence} of $f$ is defined as
        \begin{align*}
            I_w(f) := \E[\mu]{\Gamma_{w}(f)} = \E[\mu]{\norm{\abs{\nabla_w} f}^2_2}.
        \end{align*}
        Also, let $I_{w,i}(f) := \E[\mu]{\Gamma_{w,i}(f)}$ denote the influence of the $i$-th coordinate.

        \item The \textit{(Boolean) surface area} of $f$ is defined as
        \begin{align*}
            A_w(f) := \E[\mu]{\sqrt{\Gamma_wf}} = \E[\mu]{\norm{\abs{\nabla_w} f}_2}.
        \end{align*}

        \item The \textit{total squared influence} of $f$ is defined as
        \begin{align*}
            J_{w}(f) := \sum_{i=1}^n \tp{\E[\mu]{\abs{\delta_{w,(i)}} f}}^2 = \norm{\E[\mu]{|\nabla_w| f}}^2_2.
        \end{align*}
    \end{enumerate}
\end{definition}


Finally, we define the gradient estimates for Schreier graphs.

\begin{definition}[Gradient estimates]
    Let $\mu$ be a distribution supported on $X$ and $\{q_u: X \mapsto \R_+\}_{u\in U}$ be transition rates reversible for $\mu$.
    Suppose $\{w_u:\Omega \to \R_+\}_{u\in U}$ is a collection of weight functions. 
    \begin{itemize}
        \item We say that the \emph{Weak Gradient Estimate (WGE)} holds with constant $\rho_{\WGE} > 0$ if for any function $f:\Omega \to \R$ and $t\geq 0$, 
        \begin{align}\label{eq:WGE-2}
            \Gamma_w(P_t f) \leq e^{-2\rho_{\WGE} t} P_t(\Gamma_w(f)) \quad\iff\quad \norm{\abs{\nabla_w} P_t f}^2_2 \leq e^{-2\rho_{\WGE} t} P_t(\norm{\abs{\nabla_w} f}^2_2); \tag{WGE}
        \end{align}
        \item We say that the \emph{Strong Gradient Estimate (SGE)} holds with constant $\rho_{\SGE} > 0$ if for any function $f: \Omega \to \R$ and $t\geq 0$, 
        \begin{align}\label{eq:SGE-2}
            \sqrt{\Gamma_w(P_t f)} \leq e^{-\rho_{\SGE} t} P_t  \sqrt{\Gamma_w(f)} \quad\iff\quad \norm{\abs{\nabla_w} P_t f}_2 \leq e^{-\rho_{\SGE} t} P_t(\norm{\abs{\nabla_w} f}_2); \tag{SGE}
        \end{align}
        \item  We say that the \emph{Coordinate-wise Strong Gradient Estimate (CSGE)} holds with constant $\rho_{\CSGE} > 0$ if for any $\tau \in [0,\infty]$, any function $f: \Omega \to \R$, and $t\geq 0$, 
        \begin{align}\label{eq:CSGE-2}
            \norm{P_\tau \abs{\nabla_w}P_t f}_2 \leq e^{-\rho_{\CSGE} t}\norm{P_{\tau+t} \abs{\nabla_w}f}_2. \tag{CSGE}
        \end{align}
    \end{itemize}
\end{definition}


% \subsection{Basic definitions: Gammas, BE, Gradient estimates, etc.}
% Let $\Sch(G,X,U)$ denote a Schreier graph where $U$ is a generating set for $G$ that is symmetric, i.e., $U^{-1} = U$.

% Suppose that
% \begin{align*}
%     U = \bigcup_{i=1}^n U_i
% \end{align*}
% where each $U_i$ is symmetric, i.e., $U_i^{-1} = U_i$.
% We remark that $\{U_i\}$ is not necessarily a partition of $U$.


% \begin{example}\label{example:SGraph}
%     \begin{itemize}
%         \item $X = G = \Z_m^n$ and $U_i = \{k e_i: k \in [m-1] \}$. ($r=1$) 

%         \item $X = \binom{[n]}{k}$, $G = S_n$, and $U_i = \{(ij): j \in [n]\setminus\{i\}\}$. ($r=2$)
%     \end{itemize}
% \end{example}

% \bigskip

% \paragraph{Key definitions:}
% \begin{itemize}
%     \item 
%     \begin{align*}
%         \Gamma_{w,i}f(x) := \frac 12 \sum_{u \in U_i} w_u(x)q_u(x)(f(x^u) - f(x))^2.
%     \end{align*}

%     \item 
%     \begin{align*}
%         \Gamma_w f = \frac{1}{r} \sum_{i=1}^n \Gamma_{w,i} f
%         = \frac 12 \sum_{u \in U} w_u(x)q_u(x)(f(x^u) - f(x))^2.
%     \end{align*}

%     \item 
%     \begin{align*}
%         I_w(f) := \E[\mu]{\Gamma_{w}(f)}.
%     \end{align*}

%     \item 
%     \begin{align*}
%         A_w(f) := \E[\mu]{\sqrt{\Gamma_wf}}.
%     \end{align*}

%     \item 
%     \begin{align*}
%         J_{w}(f) := \frac{1}{r} \sum_{i=1}^n \tp{\E[\mu]{\sqrt{\Gamma_{w,i} f}}}^2.
%     \end{align*}

%     \item 
%         \begin{align}%\label{eq:WGE}
%             \Gamma_w P_t f \leq e^{-2\rho_{\mathrm{GE}} t} P_t \Gamma_w f \tag{WGE}
%         \end{align}
%         \item 
%         \begin{align}%\label{eq:SGE}
%             \sqrt{\Gamma_w P_t f} \leq e^{-\rho_{\mathrm{SGE}} t} P_t  \sqrt{\Gamma_w f}  \tag{SGE}
%         \end{align}
%         \item 
%         \begin{align}%\label{eq:CSGE}
%             \sum_{i=1}^n \left( P_\tau \sqrt{\Gamma_{w,i} P_t f} \right)^2 \leq e^{-2\rho_{\mathrm{CSGE}} t} \sum_{i=1}^n \left( P_{\tau+t} \sqrt{\Gamma_{w,i} f} \right)^2 \tag{CSGE($\tau$)}
%         \end{align}
% \end{itemize}

\subsection{Isoperimetric inequalities via gradient estimates}\label{sec:SGraph_iso_ineq}

In this section, we extend the isoperimetric inequalities from~\Cref{sec:Bool} to Schreier graphs. With the exception of Bobkov's inequality, the proofs do not rely on the explicit form of the operator $L_w$ and therefore carry over directly to this setting. We thus only prove Bobkov's inequality from scratch and state the remaining results without proof.

\begin{theorem}[Bobkov inequality]\label{lem:local_bobkov-sch}
    Let $\Sch(G,X,U)$ be a connected Schreier graph. Let $\mu$ be a distribution supported on $X$ and $\{q_u: X \to \R_+\}_{u \in U}$ be transition rates reversible for $\mu$. 
    Assume $\{w_u:X \to \R_+\}_{u\in U}$ is a collection of weight functions that are $b$-bounded for some $b \in (0,1]$. 
    
    If the strong gradient estimate holds with constant $\rho > 0$,
    then, for any $f:X \to [0,1]$ and $t\geq 0$, it holds
    \begin{align}\label{eq:local_bobkov-sch}
        \+I(P_t f)\leq P_t\tp{\sqrt{\+I(f)^2 + \gamma(t) \Gamma_w(f)}},
    \end{align}
    where $\gamma(t) = \frac{2(1-e^{-2\rho t})}{b \rho}$ and $\+I:[0,1] \to \R$ denotes the Gaussian isoperimetric function.

    In particular, by taking $t \to +\infty$ on both sides, it holds
    \begin{align*}
        \+I(\E[\mu]{f}) \le \E[\mu]{\sqrt{I(f)^2 + \frac{2}{b \rho} \Gamma_w(f)}}.
    \end{align*}
    % \ZC{In particular, ...(add the version with $t = \infty$)
    % }
\end{theorem}

\begin{proof}
    Fix $f : \Omega \to [0,1]$ and $t \ge 0$.
    By applying the argument to $f_\varepsilon:=\varepsilon+(1-2\varepsilon)f$, where $0<\varepsilon<1/2$, and then letting $\varepsilon\downarrow0$, we may assume that $f$ takes values in $(0,1)$.
    Following a similar argument in~\Cref{lem:local_bobkov}, it only remains to show
    \begin{align}\label{eq:dPsi-7-sch}
        L \+I^2(P_{t-s} f) - 2 \+I(P_{t-s} f) \+I'(P_{t-s} f) P_{t-s} L f+ \frac{4}{b}\Gamma_w(P_{t-s} f) - 2 \Gamma(\+I(P_{t-s} f)) \ge 0.
    \end{align}
    By definition of $L$ and $\Gamma_w$, the left-hand-side of the above inequality satisfies
    \begin{align*}
    \text{LHS of~\eqref{eq:dPsi-7-sch}}(x) &= \sum_{i=1}^n \sum_{u \in U_i} q_u(x) \tp{\+I^2(P_{t-s} f)(x^u) - \+I^2(P_{t-s} f)(x)}\\
    &- 2 \sum_{i=1}^n \sum_{u \in U_i} q_u(x) \+I(P_{t-s} f)(x) \+I'(P_{t-s} f)(x) ((P_{t-s} f)(x^u) - P_{t-s}(f)(x))\\
    &+ \frac{2}{b}\sum_{i=1}^n \sum_{u \in U_i} q_u(x) w_u(x) \tp{P_{t-s} f(x^u) - P_{t-s} f(x)}^2\\
    &-\sum_{i=1}^n \sum_{u \in U_i} q_u(x) (\+I(P_{t-s} f)(x^u) - \+I(P_{t-s} f)(x))^2,
    \end{align*}
    which is non-negative by $w_u(x) \ge b$ and
    \begin{align*}
        \+I(x)\tp{\+I(y) - \+I(x) - \+I'(x)(y-x)} + (y-x)^2 \ge 0,
    \end{align*}
    for all $x,y \in (0,1)$.
\end{proof}

\begin{theorem}[$L^1$--$L^2$ Talagrand inequality]\label{lem:l1_l2_talagrand_sch}
     Let $\Sch(G,X,U)$ be a connected Schreier graph.
     Let $\mu$ be a distribution supported on $X$ and $\{q_u: X \mapsto \R_+\}_{u \in U}$ be transition rates reversible for $\mu$.
    Assume~\eqref{eq:0-CSGE} holds with constant $\rho_{\CSGE} > 0$, and $\{w_u:\Omega \to \mathbb{R}_+\}_{u \in U}$ is a collection of weight functions that are $b$-bounded for some $b \in (0,1]$. Furthermore, suppose the log-Sobolev inequality holds for constant $\rho_{\LSI} > 0$.
    The following inequality holds for any $f:X \to \R$:
    \begin{align}\label{eq:l1_l2_talagrand_2_sch}
        \Var[\mu]{f} \leq C_1(b,\rho_{\LSI},\rho_{\CSGE}) \cdot \sum_{i=1}^n \frac{I_{w,i} (f)}{1 + \log\tp{\sqrt{I_{w,i}(f)} / \E[\mu]{\abs{\delta_{w,i}} f}}}.
    \end{align}
    A summand is interpreted as zero when $I_{w,i}(f)=0$.
    Furthermore, we have that
    \begin{align}\label{eq:l1_l2_talagrand_3_sch}
        I(f) \geq C_2(b,\rho_{\LSI},\rho_{\CSGE})\cdot \Var[\mu]{f}\tp{1 + \log \frac{I_w(f)}{J_w(f)}} .
    \end{align}
\end{theorem}

\begin{theorem}[Eldan--Gross inequality]\label{thm:eldan_gross_sch}
      Let $\Sch(G,X,U)$ be a connected Schreier graph.
     Let $\mu$ be a distribution supported on $X$ and $\{q_u: \Omega \mapsto \R_+\}_{u\in[n]}$ be transition rates reversible for $\mu$.
    Assume both~\eqref{eq:0-CSGE} and~\eqref{eq:infty-CSGE} holds with constant $\rho_{\CSGE} > 0$, and $\{w_i:X \to \mathbb{R}_+\}_{i \in [n]}$ is a collection of weight functions that are $b$-bounded for some $b \in (0,1]$. Furthermore, suppose the log-Sobolev inequality holds for constant $\rho_{\LSI} > 0$.
    For any Boolean function $f:X \to \{\pm 1\}$, it holds
   \begin{align}\label{eq:eldan_gross-sch}
        A_w(f) %\E[\mu]{\norm{\nabla_w f}_2}
        \geq C(\rho_{\LSI},\rho_{\CSGE},b) \cdot \Var[\mu]{f}\sqrt{\log\tp{1 + \frac e{J_w(f)}}}.
   \end{align}

\end{theorem}

\subsection{Establishing CSGE}
Similar to the Boolean hypercube case, we also define the canonical transition rates and weights for Schreier graphs as follows,
\begin{equation}\label{eq:canonical_SGraph}
    \forall u\in U, x\in X, q_u^\star(x) = \frac 1{2}\sqrt{\frac{\mu(x^u)}{\mu(x)}}\quad \text{and}\quad w_u^\star(x) = 2q_u^\star(x) = \sqrt{\frac{\mu(x^u)}{\mu(x)}}.
\end{equation}
Under the canonical transition rates and weights the associated semigroup is reverseible with respect to $\mu$ and in most cases, providing easy forms of gradients and weighted carr\'e du champ operators.
We introduce the following notations for the canonical transition rates and weights,
$$
    \delta^\star_u f(x) = \frac{q_u^\star(x)}{\sqrt r}(f(x^u) - f(x)), 
    \quad \grad_i^\star f(x) = (\delta^\star_u f(x))_{u\in U_i};
$$
$$
    \quad \abs{\delta^\star_i} f(x) = \norm{\grad_i^\star f(x)}_2,
    \quad \abs{\nabla^\star} f(x) = (\abs{\delta_i^\star} f(x))_{i\in[n]};
$$
$$
    \Gamma_i^\star f(x) = \Gamma_{w^\star,i} f(x)=  \sum_{u\in U_i} \delta_u^\star f(x)^2, \quad L f = \sqrt r\sum_{u\in U} \delta^\star_u f.
$$
Here we follow the notation in \Cref{sec:CSGE_boolean_hypercube} and superscript $\star$ to denote operators associated with canonical transition rates and weights.

We define $I(u) := \{i\in [n]\mid u\in U_i\}$ to be the indices of conjugacy classes to which the operator $u$ belongs and $N(S) :=  \{u\in U\mid \exists v\in S, I(u)\cap I(v) \neq \emptyset\}$ to be operators that are incident to $S$.
We interpret $I(u)$ as the set of indices that can be influenced by $u$.
Accordingly, one natural commutativity assumption about the operators is that,
\begin{equation}\label{eq:commutativity_assumption_CSGE}
    \forall u,v\in U, x\in X, I(u)\cap I(v) = \emptyset \implies x^{uv} = x^{vu},
\end{equation}
where we use $x^{uv}$ to denote $(x^u)^v$.
Under this assumption, we have the similar square property for distribution $\mu$ over Schreier graphs.
\begin{fact}\label{fact:square_property_SGraph}
    Suppose \Cref{eq:commutativity_assumption_CSGE} holds for Schreier graph $(G,X,U)$.
    For any $u,v\in U$ with $I(u)\cap I(v) = \emptyset$, we have that
    $$
        q_u^\star(x)q_v^\star(x^u) = q_v^\star(x)q_u^\star(x^v).
    $$
    holds for any $x\in X$.
\end{fact}


In this section, we generalize the sufficient condition of CSGE from boolean hyperbube to Schreier graphs.
By decomposing the relevant quantity into the local curvature term and other influence terms, we can derive a matrix-based upper bound for CSGE.
As in the Boolean hypercube setting, with the canonical transition rates and weights, the matrix can be entry-wisely upper bounded by Dobrushin influence matrix and marginal bounds.
In the more general Schreier graph setting, however, an additional lower bound on the local curvature is required.
\begin{theorem}\label{thm:SGraph_CSGE_general}
Suppose \Cref{eq:commutativity_assumption_CSGE} holds for Schreier graph $(G,X,U)$.
Let $\mu$ be a distribution on $X$ with full support. 
Let $M^\star,D^\star,D^\star_{\-{loc}}\in \R^{n \times n}$ be matrix fields satisfying that
\begin{equation}\label{eq:def_M_star_Sgraph}
    M^\star := \tp{\sup_x\norm{\tp{\frac{\ind[v\notin N(U_i)]}{2r}\tp{\sqrt{\frac{\mu(x^{uv})}{\mu(x^v)}}-\sqrt{\frac{\mu(x^u)}{\mu(x)}}}}_{u\in U_i, v\in U_j}}_2}_{i,j\in[n]};
\end{equation}
\begin{equation}\label{eq:def_D_star_Sgraph}
    D^\star := \diag\left\{\inf_{x\in X,u\in U_i} \frac {1}{2}\sum_{v\in U\setminus N(U_i)} \tp{\sqrt{\frac{\mu(x^{uv})}{\mu(x^u)}} - \sqrt{\frac{\mu(x^v)}{\mu(x)}}}\right\}_{i\in[n]}.
\end{equation}
\begin{equation}\label{eq:def_D_star_loc_Sgraph}
    D^\star_{\-{loc}} := \diag\left\{\inf_{f,x} \frac{\tp{\sqrt{\Gamma_i^\star f}L_{N(U_i)}\sqrt{\Gamma_i^\star f} - \Gamma_i^\star(f,L_{N(U_i)}f)}(x)}{\Gamma_{i}^\star f (x)} \right\}_{i\in[n]} \quad\text{where}\quad L_S(f) = \sqrt r\sum_{u\in S} \delta_u^\star f.
\end{equation}
Then, under the canonical transition rates and weights, for any $\tau \in \R_+$, \ref{eq:CSGE-2} holds with constant
\begin{align*}
    \rho =
    \lambda_{\min}\left( D_{\-{loc}}^\star + D^\star - \frac{1}{2}\left( M^\star + (M^\star)^\top \right) \right),
\end{align*}
where $\lambda_{\min}(A)$ denotes the minimum eigenvalue of a symmetric matrix $A$.
\end{theorem}

We first prove the following lemma, which provides a lower bound for the commutator-type term arising in the proof.
\begin{lemma}\label{lem:abs-commutator-Sgraph}
    Suppose \Cref{eq:commutativity_assumption_CSGE} holds for Schreier graph $(G,X,U)$.
    Let $M^\star,D^\star,D^\star_{\-{loc}}$ be matrix fields defined in \Cref{thm:SGraph_CSGE_general}.
    For any function $g: X \to \R$, it holds that for any $i\in [n]$,
    \begin{align*}
        L \abs{\delta^\star_i} g -  \frac{\inner{\grad_i^\star g}{\grad_i^\star L g}}{\abs{\delta^\star_i} g} \geq ((D^\star_{\-{loc}})_{ii}+D^\star_{ii})\abs{\delta_i^\star} g - \sum_{j\neq i} M^\star_{ij} \abs{\delta_j^\star} g 
    \end{align*}
\end{lemma}
\begin{proof}[Proof of \Cref{lem:abs-commutator-Sgraph}]
    Fix $i\in [n]$.
    Recall that the definition of $D_{\-{loc}}^\star$ in \Cref{eq:def_D_star_loc_Sgraph} derives
    \begin{equation}\label{eq:local_term_Sgraph}
        \abs{\delta_i^\star} g\cdot L_{N(U_i)} \abs{\delta_i^\star} g - \inner{\nabla_i^\star g}{\nabla^\star_i L_{N(U_i)}g} \geq (D_{\-{loc}}^\star)_{ii}\cdot(\abs{\delta_i^\star} g)^2.
    \end{equation}
    Let $Q = U\setminus N(U_i)$.
    It is sufficient to lower bound the non-local term $\abs{\delta_i^\star} g\cdot L_{Q} \abs{\delta_i^\star} g - \inner{\nabla_i^\star g}{\nabla^\star_i L_Qg}$.
    
    For notational convenience, we occasionally omit the dependence on $x$.
    Define $H_{uv}(x) = \frac{q_u^\star(x^v)}{q_u^\star(x)}$.
    By \Cref{fact:square_property_SGraph}, for any $u,v\in U$ with $I(u)\cap I(v) = \emptyset$, $H_{uv}(x) = H_{vu}(x)$ holds.
    Fix $u,v\in U$ with $I(u)\cap I(v) = \emptyset$ and $x\in X$, since $x^{uv} = x^{vu}$ and $H_{uv}(x) = H_{vu}(x)$, we have that
    \begin{align}
        (\delta_v^\star \delta_u^\star - \delta_u^\star \delta_v^\star)g &= \frac{q_v^\star q_u^\star}{r}\tp{H_{uv}(g(x^{uv}) - g(x^v)) - (g(x^u)-g(x))} \nonumber
        \\&-\frac{q_v^\star q_u^\star}{r}\tp{H_{vu}(g(x^{vu}) - g(x^u)) - (g(x^v)-g(x))} \nonumber
        \\&= \frac{q_v^\star q_u^\star}{r}(H_{uv}-1)(g(x^u) - g(x^v)) \nonumber
        \\&= \frac{H_{uv}-1}{\sqrt r}(q_v^\star \delta_u^\star - q_u^\star  \delta_v^\star) g. \label{eq:commutator-SGraph}
    \end{align}
    Let $Q = U\setminus N(U_i)$.
    Hence, we have that for any $v\in Q$,
    \begin{align}
        |\delta_i^\star| g \cdot \delta_v^\star |\delta_i^\star| g - \inner{\nabla_i^\star g}{\nabla_i^\star \delta_v^\star g}
        &= \frac{q_v^\star}{\sqrt r} (|\delta_i^\star| g\cdot  |\delta_i^\star| g(x^v) - \inner{\nabla_i^\star g}{\nabla_i^\star g(x^v)}) \nonumber
        \\&- \inner{\grad_i^\star g}{\grad_i^\star \delta_v g  + \frac{q_v^\star}{\sqrt r}(\grad_i^\star g - \grad_i^\star g(x^v))} \nonumber
        \\&\geq \sum_{u\in U_i} \delta_u^\star g \cdot \tp{\delta_v^\star \delta_u^\star - \delta_u^\star \delta_v^\star} g  \nonumber
        \\&\geq \frac 1{\sqrt r}\sum_{u\in U_i} \delta^\star_u g \cdot (H_{uv}-1)(q_v^\star \delta_u^\star - q_u^\star  \delta_v^\star) g,
    \end{align}
    where the first inequality follows from Cauchy-Schwarz inequality and the secound inequality follows from \Cref{eq:commutator-SGraph}.
    Now it is sufficient to lower bound 
    \begin{align}
        &\abs{\delta^\star_i} g\cdot L_Q \abs{\delta^\star_i} g -  \inner{\grad_i^\star g}{\grad_i^\star L_Q g} \nonumber
        \\\geq &\sum_{v\in Q}\sum_{u\in U_i} \delta^\star_u g \cdot (H_{uv}-1)\tp{q_v^\star \delta_u^\star g - q_u^\star  \delta_v^\star g} \nonumber
        \\\geq & D^\star_{ii}(\abs{\delta_i^\star} g)^2 - \frac{1}{r}\sum_{j\neq i}\sum_{v\in U_j}\sum_{u\in U_i}\ind[v\notin N(U_i)](q_u^\star(x^v)-q_u^\star) \delta^\star_u g \cdot \delta^\star_v g \nonumber
        \\\geq &D^\star_{ii}(\abs{\delta_i^\star} g)^2 - \sum_{j\neq i} M^\star_{ij} \abs{\delta_i^\star} g\cdot \abs{\delta_j^\star} g , \label{eq:non-local_term_Sgraph}
    \end{align}
    where $D^\star$ and $M^\star$ are define in \Cref{eq:def_D_star_Sgraph,eq:def_M_star_Sgraph} respectively.
    We conclude this lemma by \Cref{eq:local_term_Sgraph,eq:non-local_term_Sgraph}.
\end{proof}

Now we prove \Cref{thm:SGraph_CSGE_general}.
\begin{proof}[Proof of \Cref{thm:SGraph_CSGE_general}]
    Fix $f: X \to \R$ and $\tau, t\geq 0$.
    In the same way as in the proof of \Cref{thm:CSGE_boolean_hypercube_w=q}, it is sufficient to prove that $\frac d{ds}F_s \geq 0$ for any $s\in [0,t]$, where $F_s := e^{-2\rho s}\norm{P_{\tau+s}\abs{\nabla^\star}P_{t-s}f}_2^2$.

    Fix $s\in [0,t]$ and let $g = P_{t-s}f$.
    We have the following equation for the derivative of $F_s$,
    \begin{align}
        \frac d{ds} F_s &= 2 e^{-2\rho s} \left( - \rho \norm{P_{\tau+s} |\grad^\star| g}_2^2 
        + \inner{P_{\tau+s} |\grad^\star| g}{\frac{d}{ds} P_{\tau+s} |\grad^\star| g}
        \right) ;
        \label{eq:CSGE_potential_function_derivative-SGraph}
    \end{align}
    \begin{align}\nonumber
        \frac d{ds} P_{\tau+s} |\delta^\star_i| g
        &= P_{\tau+s} \left( L \abs{\delta^\star_i} g -  \frac{\inner{\grad_i^\star g}{\grad_i^\star L g}}{\abs{\delta^\star_i} g}\right).
    \end{align}
    By \Cref{lem:abs-commutator-Sgraph} and symmetry, we have that 
    \begin{align}
    \label{eq:derivative-CSGE-Sgraph}
        \frac d{ds} P_{\tau+s} |\delta^\star_i| g
        &\geq (D^\star_{\-{loc}}+D^\star - \frac{M^\star + (M^\star)^\top}2)P_{\tau+s}\abs{\nabla^\star}g .
    \end{align}
    Therefore, \Cref{eq:CSGE_potential_function_derivative-SGraph,eq:derivative-CSGE-Sgraph} shows that
    $$
        \frac d{ds} F_s \geq 2 e^{-2\rho s} (P_{\tau+s}\abs{\nabla^\star}g)^\top  (D^\star_{\-{loc}}+D^\star - \frac{M^\star + (M^\star)^\top}2 - \rho I)(P_{\tau+s}\abs{\nabla^\star}g) .
    $$
    Therefore, the coordinate-wise strong gradient estimate (\ref{eq:CSGE-2}) holds with constant $\rho \leq \lambda_{\min}(D_{\-{loc}}^\star+D^\star - \frac{1}{2}(M^\star+(M^\star)^\top))$ at all $x \in X$.
\end{proof}


\subsection{Application to $k$-slice}
In this section, we take the uniform distribution over $k$-slice as an example and derive isopermetric inequalites under canonical transition rates and weights through the log-concavity and \ref{eq:CSGE-2} condition.
% By \cref{thm:SGraph_CSGE_general} we have the following lemma for $k$-slice.
\begin{example}[uniform distribution on $\binom{[n]}k$]\label{example:CSGE_uniform_k_slice}
    If $\mu$ is the uniform distribution, then $M^\star = D^\star = \*0$ and $q=\frac 12, w=1$.
    Since for any $u\in U$ and $x\in \binom{[n]}k$,
    \begin{align*}
        2\sqrt 2 \cdot \delta_u^\star L f(x) = Lf(x^u) - Lf(x) &= 
        -Lf(x) + \frac 12\sum_{u^{-1}vu \in U}f(x^{uu^{-1}vu})-f(x^u)
        \\&= -Lf(x) + \frac 12\sum_{v \in U}f(x^{vu})-f(x^u) = 2\sqrt 2 \cdot L\delta_u^\star f(x),
    \end{align*}
    we conclude that
    \begin{align*}
        2\Gamma_{i}^\star(f,Lf)(x) &= \sum_{u\in U_i}\delta_u^\star f(x) \delta_u^\star Lf(x)
        = \sum_{u\in U_i}\delta_u^\star f(x) L\delta_u^\star f(x)
        \\&= \sum_{v\in U}\tp{\inner{\nabla_i^\star f(x)}{\nabla_i^\star f(x^v)} - \norm{\nabla_i^\star f(x)}_2^2}.
    \end{align*}
    Therefore, we have that for any $x\in \binom{[n]}k$ and $i\in[n]$,
    \begin{align*}
        \frac{\sqrt{\Gamma_i^\star f(x)}L\sqrt{\Gamma_i^\star f}(x) - \Gamma_i^\star(f,Lf)(x)}{\Gamma_i^\star f(x)}
        &= \sum_{v\in U}\frac{\norm{\nabla_i^\star f(x^v)}_2\norm{\nabla_i^\star f(x)}_2 - \inner{\nabla_i^\star f(x)}{\nabla_i^\star f(x^v)}}{\norm{\nabla_i^\star f(x)}_2^2}
        \\&\geq \sum_{v\in U_i}\frac{\norm{\nabla_i^\star f(x^v)}_2\norm{\nabla_i^\star f(x)}_2 - \inner{\nabla_i^\star f(x)}{\nabla_i^\star f(x^v)}}{\norm{\nabla_i^\star f(x)}_2^2}
        \\&\geq \sum_{v\in U_i}\frac{2(\delta_v^\star f(x))^2}{\norm{\nabla_i^\star f(x)}_2^2} = 2,
    \end{align*}
    where the last inequality follows from $\delta_u^\star f(x^u) = -\delta_u^\star f(x)$ and Cauchy-Schwarz inequality.
    Hence, $2\leq \lambda_{\min}(D^\star_{\-{loc}} + D^\star - \frac{M^\star+(M^\star)^\top}{2})$ and the CSGE constant is $2$.
\end{example}
\begin{theorem}[\cite{lee1998logarithmic}]
    The log-Sobolev constant of unidform distribution and Schreier graph on $\binom{[n]}{k}$ is as follows,
    $$
        \rho = \Theta\tp{\frac n{\log(1/\nu(k))}},
    $$
    where $\nu(k) = \nu(n-k) = k(n-k)/\binom{n}2$.
\end{theorem}
Plugging the above example and \ref{eq:LSI} into \cref{thm:SGraph_CSGE_general} and \cref{sec:SGraph_iso_ineq} respectively, we conclude the isoperimetric inequalities for uniform distribution on $k$ slices.
\begin{theorem}\label{thm:k_slice_isoperimetric}
    Let $\mu$ be the uniform distribution over $\binom{[n]}k$.
    Suppose the transition rates and weights are the canonical transition rates and weights we define in \eqref{eq:canonical_SGraph}.
    Then the local Bobkov inequality \eqref{eq:}, the $L^1-L^2$ Talagrand inequality \eqref{eq:}, the KKL inequality \eqref{eq:} and the Eldan-Gross inequality \eqref{eq:} holds with constant depending on $k$.
\end{theorem}


\subsection{Application to hypergrid}
% \begin{lemma}[CSGE distributions on hypergrid {$[m]^n$}]
%     Let the Scherier graph be Hamming graph we defined in \Cref{example:SGraph}.
%     Take transition rate and weight to be the canonical transition rates and weights we define in \eqref{eq:canonical_SGraph}.
%     Then CSGE condition holds ($\rho > 0$) if $\norm{K := C^{-1/2} (A+\frac{B+B^\top}{2}) C^{-1/2}}_2 < 1$ where
%     $$
%         A := \diag\left\{\sup_{x,k_1\in [Q-1]} \sum_{j\neq i}\sum_{k_2\in [Q-1]}\abs{\sqrt{\frac{\mu(x^{k_2e_j})}{\mu(x)}}-\sqrt{\frac{\mu(x^{k_1e_i,k_2e_j})}{\mu(x^{k_1e_i})}}}\right\}_{i\in[n]},
%     $$
%     $$
%         B := \tp{\ind[i\neq j]\sup_x\norm{\tp{\abs{\sqrt{\frac{\mu(x^{k_1e_i})}{\mu(x)}}-\sqrt{\frac{\mu(x^{k_1e_i,k_2e_j})}{\mu(x^{k_2e_j})}}}}_{k_1,k_2\in [m-1]}}_2}_{i,j\in[n]},
%     $$
%     $$
%         C := \diag\left\{\inf_{f,x} \frac{\sqrt{\Gamma_{w,i}f}L_{U_i}\sqrt{\Gamma_{w,i}f} - \Gamma_{w,i}(f,L_{U_i}f)}{\Gamma_{w,i} f} \right\}_{i\in[n]}.
%     $$
% \end{lemma}
% \begin{remark}
%     For each $i\in[n]$, the diagonal entry $C_{ii}$ coincides with the reinforced Bakry--Émery condition \eqref{eq:rBE} at the coordinate $i$.
% \end{remark}

Similarly in \Cref{sec:CSGE_boolean_hypercube}, we can upper bound matrix $M^\star$ and $D^\star$ by marginal bounds and Dobrushin matrix we defiened in \Cref{def:}.

\begin{lemma}\label{lem:hyper_grid_Dobrushin_to_AB_in_CSGE}
    Suppose the Scherier graph is the Hamming graph and the transition rates and weights are the canonical transition rates and weights we define in \eqref{eq:canonical_SGraph}.
    If $w_u(x)^\star \in [b,1/b]$ for $x\in X$ and $u\in U$, then
    we have that
    $$
        D^\star \geq_{\-{ew}}-\frac{(m+2b)(b^2+m-1)}{4b^3}\cdot  \diag\{R^\top \*1\}
        \qquad \text{and} \qquad M^\star \leq_{\-{ew}}\frac{(m-1)(b^2+m-1)(b+1)}{4b^3} \cdot R .
    $$
\end{lemma}
\begin{proof}
    Fix $x\in X$, $i\neq j$ and $k_1\in [m-1]$.
    WLOG, assume $x_i = 0$.
    Let $\alpha$ be the marginal distribution $\mu_j(\cdot\mid x_{-j})$, that is, $\alpha_k = \mu(X_j = k\mid X_{-j}=x_{-j})$.
    In the same way, let $\beta$ be the marginal distribution of $\mu_j(\cdot\mid x_{-j}^{k_1 e_i})$.
    The marginal lower bound is as follows,
    $$
        \forall k\in \{0\}\cup [m-1], \min\{\beta_k,\alpha_k\} \geq \frac{1}{1 + (m-1)/b^2}.
    $$
    Therefore, for any $k_2\in [m-1]$,
    \begin{align}
        2\abs{q_{k_2 e_j}^\star(x) - q_{k_2 e_j}^\star(x^{k_1 e_i})}
        &= \abs{\sqrt{\frac{\alpha_{k_2}}{\alpha_0}} - \sqrt{\frac{\beta_{k_2}}{\beta_0}}} 
        \leq \abs{\sqrt{\frac{\alpha_{k_2}}{\alpha_0}} - \sqrt{\frac{\beta_{k_2}}{\alpha_0}}} + \abs{\sqrt{\frac{\beta_{k_2}}{\alpha_0}} - \sqrt{\frac{\beta_{k_2}}{\beta_0}}} \nonumber
        \\&\leq \frac{\abs{\alpha_{k_2} - \beta_{k_2}}}{\sqrt{\alpha_{0}\alpha_{k_2}}+\sqrt{\alpha_0\beta_{k_2}}} + \sqrt{\frac{\beta_{k_2}}{\beta_0}}\frac{\abs{\alpha_0 - \beta_0}}{\alpha_0 + \sqrt{\alpha_0\beta_0}} \nonumber
        \\&\leq \frac{b^2+m-1}{2b^2}\abs{\alpha_{k_2} - \beta_{k_2}} + \frac{b^2+m-1}{2b^3}\abs{\alpha_0 - \beta_0} , \label{eq:Dobrushin_upper_bound_AB}
    \end{align}
    where the last inequality follows from the marginal bounds.
    Note that in Hamming graph, 
    $$
        \forall i\in[n], D^\star_{ii} = \inf_{x\in X, k_1\in [m-1]} \sum_{j\neq i}\sum_{k_2\in [m-1]}(q_{k_2 e_j}(x^{k_1 e_i}) - q_{k_2 e_j}(x)).
    $$
    Therefore, given $x,k_1$, we have that by \Cref{eq:Dobrushin_upper_bound_AB},
    \begin{align*}
        \sum_{j\neq i}\sum_{k_2}\abs{q_{k_2 e_j}(x) - q_{k_2 e_j}(x^{k_1 e_i})}
        \leq \frac 12\sum_{j\neq i} \tp{\frac{b^2+m-1}{b^2} + \frac{m(b^2+m-1)}{2b^3}}\TV{\mu(\cdot\mid x_{-j})}{\mu(\cdot\mid x_{-j}^{k_1 e_i})}.
    \end{align*}
    This implies that $D^\star \geq_{\-{ew}} -\frac{(m+2b)(b^2+m-1)}{4b^3}\diag\{R^\top \*1\}$.
    On the other hand, \Cref{eq:Dobrushin_upper_bound_AB} shows that
    for any $x \in X$,
    $$
        \tp{\abs{\sqrt{\frac{\mu(x^{k_1e_i})}{\mu(x)}}-\sqrt{\frac{\mu(x^{k_1e_i,k_2e_j})}{\mu(x^{k_2e_j})}}}}_{k_1,k_2\in [m-1]} \leq_{\-{ew}} \frac{(b^2+m-1)(b+1)}{2b^3}\cdot R_{ij} J,
    $$
    where $J$ is the all one matrix.
    Hence, $M^\star \leq_{\-{ew}} \frac{(m-1)(b^2+m-1)(b+1)}{4b^3} \cdot R$.
\end{proof}

\begin{lemma}[Local curvature lower bound]\label{lem:local_curvature}
    Suppose the Scherier graph is the Hamming graph and the transition rates and weights are the canonical transition rates and weights we define in \eqref{eq:canonical_SGraph}.
    If $w_u(x)^\star \in [b,1/b]$ for $x\in X$ and $u\in U$, then we have that
    $$
        \inf_{f,x} \frac{\sqrt{\Gamma_{w,i}f}L_{U_i}\sqrt{\Gamma_{w,i}f} - \Gamma_{w,i}(f,L_{U_i}f)}{\Gamma_{w,i} f} \geq b(\sqrt{m-1}+1)-(m-1)(\frac 1b-b).
    $$
\end{lemma}
\begin{proof}
    Fix $i\in[n]$ and $f\in \R^\Omega$.
    We use $q_{k} = q_{k e_i}$ to simplify the notation.
    Also, we omit $(x)$ for simplicity.
    Let $d_k = f(x^{k e_i}) - f(x)$.
    For the first part $\sqrt{\Gamma_{w,i} f}L_{U_i}\sqrt{\Gamma_{w,i} f}$, we have the following,
    \begin{align}
        \sum_{k=1}^{m-1} q_k\sqrt{\frac{\Gamma_{w,i} f(x^k)}{\Gamma_{w,i}f}} 
        &= \sum_{k=1}^{m-1}q_k\sqrt{\frac{\sum_{\ell = 1}^{m-1}q_\ell(x^k)^2(f(x^{(k+\ell)e_i}) - f(x^{k e_i}))^2}{\sum_{\ell=1}^{m-1}q_\ell^2d_\ell^2}} \nonumber
        \\& = \sum_{k=1}^{m-1}\sqrt{\tp{\sum_{\ell = 1}^{m}q_\ell^2(d_\ell - d_k)^2} / \tp{\sum_{\ell=1}^{m-1}q_\ell^2d_\ell^2}} &&(\text{Rearrange}) \nonumber
        \\&\geq \sqrt{\tp{\sum_{\ell = 1}^m q_\ell^2\tp{\sum_{k=1}^{m-1} d_\ell - d_k}^2 }/ \tp{\sum_{\ell=1}^{m-1}q_\ell^2d_\ell^2}} &&(\text{Triangle ineq.}) \nonumber
        \\&\geq b\sqrt{\tp{\sum_{\ell = 1}^m \tp{\sum_{k=1}^{m-1} d_\ell - d_k}^2 }/ \tp{\sum_{\ell=1}^{m-1}d_\ell^2}} &&(\forall \ell,k,q_\ell/q_k \in [b,1/b]) \nonumber
        \\&= b\sqrt{(m-1)^2 - (m-2)\tp{\sum_{k=1}^{m-1}d_k}^2/\tp{\sum_{k=1}^{m-1}d_k^2}} \nonumber
        \\&\geq b\sqrt{(m-1)}. &&(\text{AM-GM}) \label{eq:Local_curvature_part1}
    \end{align}
    Define $Q = \sum_{k=1}^{m-1} q_k$.
    Therefore, 
    $$
        \frac{\sqrt{\Gamma_{w,i} f}L_{U_i}\sqrt{\Gamma_{w,i} f}}{\Gamma_{w,i} f} \geq b\sqrt{m-1} - Q.
    $$
    For the second part $-\Gamma_{w,i}(f,L_{U_i} f)$, it is lower bounded as follows,
    \begin{align}
        -2\Gamma_{w,i}(f,L_{U_i} f)
        &= -\sum_{k=1}^{m-1} q_k^2 d_k \sum_{\ell=1}^{m-1} \tp{q_l(x^k)(f(x^{(k+\ell)e_i}) - f(x^k)) - q_\ell d_\ell} \nonumber
        \\&=\sum_{k=1}^{m-1}\sum_{\ell=1}^{m-1}q_k^2q_\ell d_kd_\ell
        + \sum_{k=1}^{m-1}\sum_{\ell=1}^m q_kq_\ell d_k^2
        - \sum_{k=1}^{m-1}\sum_{\ell=1}^m q_kq_\ell d_k d_\ell \nonumber
        \\&= (qd)^\top (q\*1^\top + (Q+1)\cdot \diag\{q\}^{-1} - \*1\*1^\top)  (qd) \nonumber
        \\&\geq \norm{qd}_2^2 \lambda_{\min}\tp{\frac{(q-\*1)\*1^\top + \*1(q-\*1)^\top}{2} + (Q+1)\cdot \diag\{q\}^{-1}} \label{eq:Local_curvature_part2} .
    \end{align}
    Since $\lambda_{\min}(A+B)\geq \lambda_{\min}(A)+\lambda_{\min}(B)$ holds for symmetric matrices $A,B$, we have that,
    $$
        \lambda_{\min}\tp{\frac{(q-\*1)\*1^\top + \*1(q-\*1)^\top}{2} + (Q+1)\cdot \diag\{q\}^{-1}} \geq \frac{\inner{q-\*1}{\*1}-\sqrt{m-1}\norm{q-\*1}_2}{2} + b(Q+1) .
    $$
    Since $2\Gamma_{w,i} f = \norm{qd}_2^2$, by \Cref{eq:Local_curvature_part1} and \Cref{eq:Local_curvature_part2}, we conclude that,
    \begin{align*}
        &\frac{\sqrt{\Gamma_{w,i} f}L_{U_i}\sqrt{\Gamma_{w,i} f} - \Gamma_{w,i}(f,L_{U_i} f)}{\Gamma_{w,i} f}
        \\\geq& b\sqrt{m-1} - Q +  \frac{Q-(m-1)-\sqrt{m-1}\norm{q-\*1}_2}{2} + b(Q+1)
        \\=& b\sqrt{m-1} -\frac{m-1-2b}{2} + (b-\frac{1}{2})Q - \frac{\sqrt{m-1}\norm{q-\*1}_2}{2} 
        \\\geq& b\sqrt{m-1} -\frac{m-1-2b}{2} + (m-1)\min\{b^2-\frac{b}{2},1-\frac{1}{2b}\} - \frac{(m-1)(1-b)}{2b} 
        \\\geq& b\sqrt{m-1} -\frac{m-1-2b}{2} + (m-1)(b-\frac 1{2b}) - \frac{(m-1)(1-b)}{2b} 
        \\=& b(\sqrt{m-1}+1)-(m-1)(\frac 1b-b),
    \end{align*}
    \ZC{Formatting.}
    where the second inequality follows from $Q \in [(m-1)b,\frac {m-1}b]$ and $\norm{q}_2 \leq \sqrt{m-1}(1-b)/b$.
\end{proof}


\Cref{lem:hyper_grid_Dobrushin_to_AB_in_CSGE} and \Cref{lem:local_curvature} immediately gives the upper bound of $\norm{K}_2$.
\begin{lemma}\label{lem:hyper_grid_Dobrushin_to_CSGE_K}
    Let $w = q$ with $q_u(x) = \sqrt{\mu(x^u)/\mu(x)} \in [b,1/b]$ for any $x\in [m]^n$ and $i\in [n]$. 
    If $m\geq 2$ and $b>\sqrt{\frac{m-1}{m+\sqrt{m-1}}}$, we have that $\norm{K}_2 \leq C_{cons}\cdot (\norm{D}_1 + \norm{D}_2)$, where
    \begin{equation}\label{eq:hyper_grid_CSGE_C_cons}
        C_{cons} = \frac{m(b^2+m-1)(b+1)}{2b^2(b^2(m+\sqrt{m-1})-m+1)} .
    \end{equation}
\end{lemma}
\begin{proof}
    By \Cref{lem:local_curvature}, $C \geq_{\-{ew}} b(\sqrt{m-1}+1)-(m-1)(\frac 1b-b) I$.
    Therefore, by \Cref{lem:hyper_grid_Dobrushin_to_AB_in_CSGE}, 
    we have that
    $$
        K \leq_{\-{ew}} \tp{b(m+\sqrt{m-1}) - \frac{m-1}{b}}^{-1}\frac{m(b^2+m-1)(b+1)}{2b^3}\tp{\diag\{D^\top \*1\} + \frac{D+D^\top}{2}}.
    $$
    Hence, $\norm{K}_2 \leq C_{cons}\cdot(\norm{D}_1 + \norm{D}_2)$.
\end{proof}
\begin{corollary}
    Let $w = q$ with $q_u(x) = \sqrt{\mu(x^u)/\mu(x)} \in [b,1/b]$ for any $x\in [m]^n$ and $i\in [n]$. 
    If $m\geq 2$ and $b>\sqrt{\frac{m-1}{m+\sqrt{m-1}}}$, the coordinate-wise strong gradient estimate \eqref{eq:CSGE} holds with constant $\rho \leq b(m+\sqrt{m-1})-(m-1)/b-\frac{m(b^2+m-1)(b+1)}{2b^3}\tp{\norm{D}_1+\norm{D}_2}$.
\end{corollary}

Recall that \Cref{thm:marton_Dobrushin_to_LSI} gives an \ref{eq:LSI} constant for the Gibbs sampler in terms of the Dobrushin matrix.
For hypergrid, the discrepancy between the Gibbs generator and our generator can be bounded by one-site marginal bounds.
Therefore, together with the upper bound on $\norm{K}_2$, we conclude the following theorem: under suitable Dobrushin conditions and marginal bounds, the isopermetric inequalities in \ref{sec:SGraph_iso_ineq} hold.
\begin{theorem}\label{thm:hyper_grid_Dobrushin_to_iso}
    Consider the Schreier Graph defined in the first item in \Cref{example:SGraph}.
    Let $\mu$ be a distribution on $[m]^n$ with full support.
    Take the transition rate to be $q_i(x) = \sqrt{\mu(x^i)/\mu(x)}$ and set $w = q$.
    Suppose that for any $x\in [m]^n, i\in[n]$, $q_i(x)\in [b,1/b]$ with 
    \begin{equation}\label{eq:hyper_grid_iso_cond}
        b > \sqrt{\frac{m-1}{m+\sqrt{m-1}}}\quad \text{and}\quad C_{cons}\cdot(\norm{D}_2 + \norm{D}_1) < 1,
    \end{equation}
    where $C_{cons}$ is defined in \Cref{eq:hyper_grid_CSGE_C_cons} .
    Then the local Bobkov inequality \eqref{eq:}, the $L^1-L^2$ Talagrand inequality \eqref{eq:}, the KKL inequality \eqref{eq:} and the Eldan-Gross inequality \eqref{eq:} holds with constant depending on $\norm{D}_1$, $\norm{D}_2$ and $b$.
\end{theorem}
\begin{proof}
    Since $C_{cons}(\norm{D}_2 + \norm{D}_1) < 1$, $C_{cons}\geq \frac{m(m-1)}{2(1+\sqrt{m-1})} > 1$ for $m\geq 3$ and $C_{cons}\geq 2$ for $m=2$, $\norm{D}_2 < 1$ holds.
    By \Cref{thm:marton_Dobrushin_to_LSI}, since
    $$
        \inner{f}{(I-G)f}_\mu = \E[\mu]{\frac 1{2n}\sum_{i=1}^n  \sum_{k=1}^{m-1} \frac{q_{k e_i}(x)^2}{1 + \sum_{\ell=1}^{m-1}q_{\ell e_i}(x)^2}\tp{f(x^{k e_i}) - f(x)}^2}
        \leq \frac 1{nmb^2}\E[\mu]{\Gamma(f,f)},
    $$
    the \ref{eq:LSI} holds with constant $\frac{mb^4(1-\norm{D}_2)^2}{m-1+b^2}$.
    Therefore, both CSGE condition and hypercontractivity holds.
    By \Cref{sec:SGraph_iso_ineq}, all the isopermetric inequalities hold.
\end{proof}

\subsubsection{Zero-field Potts model}
As an application, let us consider the zero-field Potts model.
A distribution $\mu$ over $[m]^n$ is a Potts model with interaction matrix $J\in \R^{n\times n}$ if
$$
    \forall x\in [m]^n, \mu(x)\propto \exp\tp{\sum_{1\leq i<j\leq n}J_{ij}\ind[x_i = x_j]} ,
$$
where $J$ is a symmetric matrix and ahs zero diagonal.
The interaction between each pair of sites depends only on whether their spins agree.

In the following theorem, we establish the isopermetric inequalities in \Cref{sec:SGraph_iso_ineq} under an explicit smallness condition on the induced $1$-norm of the interaction matrix.
\begin{theorem}[zero-field Potts model]\label{thm:Potts}
    Suppose that $\mu$ is an zero-field Potts model parameterized by an interaction matrix $J$ satisfying 
    \begin{equation}\label{eq:Potts_cond}
        t < \log\tp{1+\frac{1+\sqrt{m-1}}{m-1}}
        \quad \text{and} \quad \frac{tm(e^{-t}+m-1)(e^{-t/2}+1)}{2e^{-t}(e^{-t}(m+\sqrt{m-1})-m+1)} < 1,
    \end{equation}
    where $t := \norm{J}_1$.
    Take the transition rate to be $q_i(x) = \sqrt{\mu(x^i)/\mu(x)}$ and set $w = q$.
    Then the local Bobkov inequality \eqref{eq:local_bobkov-sch}, the $L^1$--$L^2$ Talagrand inequality \eqref{eq:}, the KKL inequality \eqref{eq:l1_l2_talagrand_2_sch} and the Eldan--Gross inequality \eqref{eq:eldan_gross-sch} holds with constant depending on $\norm{J}_1$.
\end{theorem}
\begin{proof}
    By \Cref{thm:hyper_grid_Dobrushin_to_iso}, it is sufficient to check \Cref{eq:hyper_grid_iso_cond}.
    Since for any $x\in [m]^n$ and $ke_i\in U$,
    $$
        \sqrt{\frac{\mu(x^{ke_i})}{\mu(x)}}
        = \exp\tp{\frac 12\sum_{j\neq i} J_{ij}(-\ind[x_j=x_i] + \ind[x_j\equiv x_i+k \mod m]) }
        \in \left[e^{-t/2},e^{t/2}\right] ,
    $$
    we can set the marginal bound $b = e^{-t/2}$.
    
    Fix $x\in [m]^n$ and $ke_j\in U$, define 
    $$
        \varphi(a) = \frac{\mu_i(x_i=a\mid x_{-i}^{k e_j})}{\mu(x_i=a\mid x_{-i})} = \exp\tp{J_{ij}(\ind[x_j+k=a]-\ind[x_j=a])}.
    $$
    Therefore, for any $S\in 2^m$,
    $$
        \frac{\mu_i(x_i\in S\mid x_{-i}^{k e_j})}{\mu(x_i\notin S\mid x_{-i}^{k e_j})}
        \leq e^{2\abs{J_{ij}}}\frac{\mu_i(x_i\in S\mid x_{-i})}{\mu(x_i\notin S\mid x_{-i})}
    $$
    $$
        \implies \mu_i(x_i\in S\mid x_{-i}^{k e_j}) - \mu_i(x_i\in S\mid x_{-i})
        \leq \frac{e^{2\abs{J_{ij}}}}{e^{2\abs{J_{ij}}}-1+1/\mu_i(x_i\in S\mid x_{-i})} - \mu_i(x_i\in S\mid x_{-i}) .
    $$
    Since given $a > 1$, $\sup_{b\in [0,1]}\frac{a}{a-1+1/b}-b = \frac{\sqrt a - 1}{\sqrt a + 1}$, $\mu_i(x_i\in S\mid x_{-i}^{k e_j}) - \mu_i(x_i\in S\mid x_{-i}) \leq \tanh(\frac {\abs{J_{ij}}}2)$.
    Therefore,
    $$
        \TV{\mu_i(\mid x_{-i}^{k e_j})}{\mu_i(\mid x_{-i})} = \sup_S \mu_i(i\in S\mid x_{-i}^{k e_j}) - \mu_i(i\in S\mid x_{-i}) \leq \tanh(\frac{\abs{J_{ij}}}{2}) .
    $$
    We have that $D \leq_{\-{ew}} \tanh(\abs{J/2}) \leq_{\-{ew}} \abs{J}/2$ and $\max\{\norm{D}_1,\norm{D}_2\} \leq t/2$.
    Therefore, \Cref{eq:hyper_grid_iso_cond} is implied by \Cref{eq:Potts_cond}.
\end{proof}
\begin{remark}
    For some specific numbers, the resulting numerical thresholds are summarized in the following table.
    \ZC{Remove it? Too AI.}
    \begin{table}[htbp]
    \centering
    \begin{tabular}{c|ccc}
        \hline
        $m$ & Positivity threshold & CSGE threshold $\rho_{CSGE}>0$
            & Isoperimetric threshold \\
        \hline
        $2$ & $1.098612$ & $0.287314$ & $\norm{J}_1<0.287314$ \\
        $3$ & $0.791683$ & $0.175290$ & $\norm{J}_1<0.175290$ \\
        $4$ & $0.647461$ & $0.121790$ & $\norm{J}_1<0.121790$ \\
        $5$ & $0.559616$ & $0.090917$ & $\norm{J}_1<0.090917$ \\
        $6$ & $0.499085$ & $0.071131$ & $\norm{J}_1<0.071131$ \\
        \hline
    \end{tabular}
    \caption{Numerical sufficient conditions for the zero-field Potts model.}
    \label{tab:potts-thresholds}
\end{table}
\end{remark}